\chapter{Related Work}
\label{ch:related}

The work builds on systems and papers from a field that is moving quickly,
and they are grouped here by the part of the stack that each one addresses.
Earlier chapters cite most of them where they are used; here they are
compared.

\section{Disaggregated Inference}

DistServe \citep{zhong2024distserve} and Splitwise \citep{patel2024splitwise}
established that prefill and decode should be separated, and showed the gain
in service quality per unit of hardware. Both assume a data-center network
and treat the cache transfer as one bulk copy whose cost is bandwidth.
Mooncake \citep{qin2024mooncake} makes the cache the center of the
architecture and adds storage for it. NVIDIA's Dynamo \citep{dynamo} and the
llm-d project \citep{llmd} package the pattern for clusters, and LMCache
\citep{lmcache2025,lmcache_blog} provides a cache layer for the same serving
engine used here. We take the split as given and study one aspect that those systems do not: the mix of traffic on the link between the two halves,
on hardware without the data-center network they assume.

\section{Libraries for Cache Transfer}

NIXL \citep{nixl} is the transfer library under the serving engine's standard
connector, with back ends for UCX and libfabric. It is the baseline of E8.
fabric-lib \citep{licker2026fabriclib} is the closest work in spirit: a
point-to-point RDMA library written for language-model systems, which makes
the choice of completion signal part of the design; the counter mechanism in
\cref{tab:signals} follows it. UCCL \citep{uccl} is a communication library
for GPUs that covers collectives and point-to-point transfers including cache
transfer. None of the three separates traffic by urgency, and none targets
software RoCE. OnePiece \citep{chen2026onepiece} is a large-scale inference
system built on RDMA for a different workload, at a scale that cannot be tested here. Ultra Ethernet, the industry's designed successor to RoCE, is described
in \cref{sec:uec} and not here, because it adopts the interface this
middleware is written to and so bears on the design directly.

\section{RDMA over Ethernet at Scale}

Running RDMA over Ethernet in large deployments is a mature subject. DCQCN
\citep{zhu2015dcqcn} is the congestion control used with RoCEv2, later work
questioned whether RoCE needs a lossless network at all
\citep{mittal2018irn}, and an account of RoCE for training clusters at very
large scale describes the engineering involved \citep{gangidi2024meta}. Work on the control plane shows that setting connections up can dominate in
elastic settings \citep{zhang2025swift}. High-performance computing reached
Ethernet from the other side: Slingshot \citep{desensi2020slingshot} is an
Ethernet-compatible interconnect with congestion control and traffic classes
of its own, it is programmed through libfabric \citep{libfabric_repo}, and
its traffic classes are the in-network counterpart of the separation that
the middleware performs at the endpoints. All of this concerns many machines
and switches. The present work has two machines and a cable, where congestion
control is not the issue and the interesting effects are inside one host.

\section{Transport Interfaces and Their Cost}

UCX \citep{shamis2015ucx} and libfabric \citep{grun2015ofi} are the two
portable interfaces over RDMA. A per-operation study of UCX over InfiniBand
\citep{papadopoulou2017ucx} supplied the instruction-counting method used in
E5, and the two have been compared across applications on a production
supercomputer \citep{bareford2022}.
An engineering report on a memory leak caused by a transport library's hooks
in a production split-serving system \citep{mistral2026leak} is a reminder
that the interface chosen brings its side effects with it. We measure one narrow quantity in this space: the cost of libfabric relative to raw verbs for small messages, with everything else held equal.

\section{Small Messages and Collectives}

Communication libraries for multi-GPU computation face the same small-message
problem from another side. NCCL \citep{nccl} is the standard library for
collectives, TACCL \citep{shah2023taccl} synthesizes collective algorithms for
a topology, and recent work reports that for small collectives the time is
dominated by software overhead and signaling, and uses flags and sentinel
buffers polled in memory to avoid completion queues \citep{shen2026sol}. The
flag and sentinel mechanisms of \cref{tab:signals} come from that line of
work. Scheduling at the level of whole jobs, as in Cassini
\citep{rajasekaran2024cassini}, is complementary: it decides when jobs use
the network, whereas the middleware decides which message goes first.

\section{RDMA Without Dedicated Hardware}

Soft-RoCE \citep{ibta_softroce2015} is what makes the lab pair possible. An
emulated RDMA network card for virtual machines \citep{rocm_ernic} and a
library for RDMA between consumer GPU machines without a managed switch
\citep{libmesh} address neighboring cases, and three modules were ported
from the second (\cref{sec:reused}). Tesla's TTPoE \citep{ttpoe} is a
different answer to the same need: a transport protocol that runs directly
over Ethernet, published with a reference kernel implementation. Its kernel
modules were built on the testbed, and it is not part of the comparison.

\section{Compressing and Storing the Cache}

Cache compression is active. KIVI \citep{liu2024kivi} quantizes the cache in
place on the GPU; transform coding \citep{staniszewski2026kvtc} targets compact storage, and
low-rank compression \citep{zhang2024lorc} reduces the size of the cache in
memory. TurboQuant
\citep{zandieh2025turboquant} and QJL \citep{zandieh2024qjl} are the methods
the quantizer here is assembled from, and a storage vendor has combined
TurboQuant with flash storage for cache restore
\citep{alvarez2026flashblade}. No new compression method is proposed here. The claim is narrower: a packed format with measured sizes, applied per
transfer under a policy.

For storage, the hierarchy of GPU memory, host memory, local flash and a
shared file system is described by several vendors
\citep{hsu2026fourtier,hsu2026kv,ibm2025storage}. The storage tier here
implements one level of it on one machine.

\section{Rule-Based Routing}

The rule tree of \cref{sec:policy} borrows its form from request routers
for the same serving engine, which combine signals with rules to choose a
model for each request \citep{liu2026router,wang2025router}. Here the same
form chooses a path for each transfer.

\section{Position of This Work}

Taken together: the existing systems move the cache as one kind of traffic,
over hardware built for it. The present work treats the traffic as mixed, on
commodity Ethernet with a software RDMA device and no GPU-direct path, and
asks whether classification and scheduling help there.
\waiting{results from E3 and E8: add one sentence that places the measured
result among these systems}
